\scanpage{15}
\handwritten{\mspage{6}}
Soit $P^m(X)$ la partie primitive de la cohomologie $H^m(X,\underline{Q}_{\unsure{\chi|\ell}})$, i.e. le noyau du composé
\[
H^m(X)\handwritten{\longrightarrow}H^m(\overset{\handwritten{\sim}}{X})\handwritten{\longrightarrow}H^m(Y)\handwritten{\longrightarrow}H^0(S,R^m f_{\handwritten{*}}(\underline{Q}_{\chi}))\ .
\]
On trouve donc grâce à I) un homomorphisme (dit ``de Griffiths'')
\[
P^m(X)\handwritten{\longrightarrow}H^1(S,R^{m-1}f_{\handwritten{*}}(\underline{Q}_{\chi}))
\]
(en fait, on y pourrait remplacer $S$ par $P^1_k$, mais \deleted{\unsure{nous at|nous et}} peu\unsure{.} importe ici)
D'ailleurs \deleted{\unsure{l'hom|lthom}} on a un homomorphisme d'immersion canonique
% Objects, labels and comma are typed; the three arrows are pen additions.
\diagram{\begin{tikzcd}[ampersand replacement=\&,column sep=small,row sep=small]
Y\handwritten{\arrow[rr]}\arrow[rr,phantom,"i"]\handwritten{\arrow[dr]}\arrow[dr,phantom,"f"']\&\&X_S\handwritten{\arrow[dl]}\arrow[dl,phantom,"g"]\\
\&S\&{,}
\end{tikzcd}}
induisant
\[
i^{\handwritten{*}}:\quad H^1(S,R^{m-1}g_{\handwritten{*}}(\underline{Q}_{\chi}))=H^1(S,H^{m-1}(X,\underline{Q}_{\chi}))\handwritten{\longrightarrow}H^1(S,R^{m-1}f_{\handwritten{*}}(\underline{Q}_{\chi})).
\]
Ceci dit, \uline{l'image inverse par l'homomorphisme de Griffiths plus haut de} $\mathrm{Im}\ i^{\handwritten{*}}$ \uline{est réduite à zéro}.

Ceci se fait par des calculs explicites essentiellement triviaux, compte tenu de la structure cohomologique connue d'une variété éclatée.

5.\quad Mettant ensemble 3) et 4), on trouve ceci:

% The interlinear insertion is typed; a pen guide brings it before "soit".
\deleted{Sous} Avec les notations de 4), \inserted{supposant $m=2n$,} soit $\unsure{s|\deleted{t}s|\deleted{s}t}$ le point générique de $P^1_k$, et \uline{supposons que l'immersion} $Y_{\bar{\unsure{s|\deleted{t}s|\deleted{s}t}}}\handwritten{\longrightarrow}X\mathbin{\ooalign{\hfil$\mathrm{x}$\hfil\cr\hfil$\unsure{\mathrm{O}|\handwritten{\bigcirc}}$\hfil\cr}}_k k(\bar s)$ \uline{de la section hyperplane} \unsure{\uline{générique}|g\uline{énérique}} \deleted{dans} s\uline{atisfasse aux conditions} a) \uline{et} b) \uline{de} 2).\deleted{(Alors pour On suppose de plus $m=2n$} \uline{Alors pour tout} $z\handwritten{\in}Z^n(X)$ \uline{dont la restriction à} $Y_{\bar s}$ \uline{est algébriquement équivalent à zéro}, $z$ \uline{est cohomologiquement équivalent à zéro} (mod torsion). \deleted{En d'autr} Par suite, le groupe quotient du groupe des cycles cohomologiquement équivalents à zéro \deleted{mod torsion} sur $Y_{\bar s}$, modulo ceux qui sont algébriquement équivalents à zéro (groupe qui mérite le nom de \uline{groupe de Griffiths} $\mathrm{Grif}\unsure{\handwritten{(}|t}Y_{\bar s})$ de $Y_{\bar s_{\deleted{x}}}$),\deleted{contient} \unsure{tensorisé} par $\underline{Q}$, contient comme sous-espace le sous-espace $P^{2n}_{\mathrm{alg}}(X)$ formé des classes de cohomologie primitives sur $X$ qui sont \uline{algébriques}, - vectoriel qui a
% archive sheet number: 14
